\section*{Chapter \ref{ch:qm:bayesian-methods} --- Bayesian Methods}

\begin{solution}{exo:qm:bayesian-methods:1}
$\mathrm{Beta}(80, 70)$, mean $80/150 = 0.533$. With the flat prior, $\mathrm{Beta}(31, 21)$, mean $31/52 = 0.596$: the informative prior
pulls the 60\% hit rate most of the way back to a coin.
\end{solution}

\begin{solution}{exo:qm:bayesian-methods:2}
$\sigma^2 = 0.2$, $\tau^2 = 0.16$, $B = 0.2/0.36 = 0.556$. Posterior mean $0.5 + 0.444 \times 1.0 = 0.94$; posterior variance $0.2 \times
0.16/0.36 = 0.089$, standard deviation 0.30.
\end{solution}

\begin{solution}{exo:qm:bayesian-methods:3}
$n(1 - \rho)/(1 + \rho) = 10\,000 \times 0.2/1.8 = 1\,111$.
\end{solution}

\begin{solution}{exo:qm:bayesian-methods:4}
$B = \sigma^2/(\sigma^2 + \hat\tau^2) = \sigma^2/s_x^2$, so the estimate keeps $1 - \sigma^2/s_x^2$ of each deviation. James--Stein toward the
grand mean keeps $1 - (k - 3)\sigma^2/\sum(x_i - \bar x)^2 = 1 - \frac{k-3}{k-1}\sigma^2/s_x^2$: the same \omterm{def:qm:estimation:estimator}{estimator} with a slightly smaller
shrinkage, the difference vanishing as $k$ grows.
\end{solution}

\begin{solution}{exo:qm:bayesian-methods:5}
$B = \sigma^2/(\sigma^2 + \tau^2)$ with $\sigma^2 = 1/T$ is one half when $1/T = \tau^2$: with $\hat\tau^2 = 0.168$, $T = 6.0$ years (6.25 with the
true $\tau = 0.4$). Six years of record are needed before a manager's own history counts as much as the platform's prior.
\end{solution}

\begin{solution}{exo:qm:bayesian-methods:6}
Given $\theta$ and $m$, the likelihood of $\tau^2$ is $\prod_i\tau^{-1}e^{-(\theta_i - m)^2/(2\tau^2)} = (\tau^2)^{-k/2}e^{-S/(2\tau^2)}$ with $S =
\sum_i(\theta_i - m)^2$. A flat prior on $\tau$ is $p(\tau^2) \propto (\tau^2)^{-1/2}$, so the conditional is proportional to
$(\tau^2)^{-(k-1)/2 - 1}e^{-S/(2\tau^2)}$: inverse gamma with shape $(k - 1)/2$ and scale $S/2$.
\end{solution}

\begin{solution}{exo:qm:bayesian-methods:7}
\texttt{best\_manager\_average()}: the best record averages 2.09, the chosen manager's true \omterm{def:qm:estimation:sharpe}{Sharpe ratio} 1.01, the empirical-Bayes estimate
1.01 and the next three years 1.02. Selection adds 1.08 to the record; the shrunk estimate removes it on average.
\end{solution}

\begin{solution}{exo:qm:bayesian-methods:8}
Shrinking toward zero with the James--Stein factor treats the signals as draws from one population centred at zero, and it does reduce the
average error of the estimates. It does not undo the search: the factor depends on the spread of the reported \omterm{def:qm:estimation:sharpe}{Sharpe ratios}, and if only the
survivors of a larger search are reported, that spread and the target are both wrong, so the selected signals remain overstated. The prior
must be fitted to everything tried, and the tests of chapter~12 still apply.
\end{solution}

\begin{solution}{pb:qm:bayesian-methods:1}
\textbf{1.} Best 2.41; mean 0.51, standard deviation 0.71.
\textbf{2.} Twelve.
\textbf{3.} $1/\sqrt3 = 0.58$: the records spread more than noise alone explains.
\textbf{4.} $\hat m = 0.51$, $\hat\tau = \sqrt{0.71^2 - 1/3} = 0.41$.
\textbf{5.} $B = 0.333/(0.333 + 0.168) = 0.665$.
\textbf{6.} $0.51 + 0.335 \times 1.90 = 1.14$, standard deviation 0.33.
\textbf{7.} 1.20 (factor 0.36 instead of 0.335).
\textbf{8.} 1.82 and 1.45 become 0.95 and 0.83; the ranking is unchanged, since the shrinkage is the same for everyone.
\textbf{9.} True ratios 1.07, 0.96 and 1.68; next three years 1.69, 0.57 and 0.99. The best record was not the best manager.
\textbf{10.} Predictive $\mathcal N(1.14, 0.67^2)$: $\Phi(-1.14/0.67) = 4.3\%$.
\textbf{11.} From 0.335 to 0.149, a cut of 55\%.
\textbf{12.} Against the truth, 0.336 to 0.121 (64\%); against the next three years, 0.672 to 0.457 (32\%).
\textbf{13.} Best record 2.09, true ratio 1.01, shrunk estimate 1.01 (next three years 1.02).
\textbf{14.} Posterior mean 1.12, 95\% \omterm{def:qm:bayesian-methods:predictive}{credible interval} $(0.41, 2.01)$: wider than the plug-in interval, because $\tau$ is uncertain.
\textbf{15.} Mean 0.40, 95\% interval $(0.13, 0.66)$, with a 7\% chance that $\tau < 0.2$, in which case the managers barely differ and the
shrinkage should be much stronger.
\textbf{16.} Not on this evidence alone: the expected \omterm{def:qm:estimation:sharpe}{Sharpe ratio} is 1.14, very good but not 2.4, and the capital should be sized to 1.14
with its uncertainty, with more added as the record lengthens.
\textbf{17.} A prior fitted to comparable external managers, wider (larger $\tau$) and probably centred lower, since hiring selects on
past records too.
\textbf{18.} When true \omterm{def:qm:estimation:sharpe}{Sharpe ratios} have heavy tails, a few genuinely exceptional managers exist and the normal prior pulls them back as hard
as the lucky ones; a Student-$t$ prior shrinks large deviations less.
\textbf{19.} Named result: \emph{the allocator's shortlist}: the best of fifty three-year records, 2.41, has an empirical-Bayes posterior
mean of 1.14; shrinkage cuts the error against the next three years' records by 31\% on this platform and 32\% on average.
\textbf{20.} Because the best of many noisy records is selected partly for its noise, and the prior removes the part of the deviation that
noise explains.
\end{solution}

\begin{solution}{iq:qm:bayesian-methods:1}
A 95\% confidence interval is produced by a procedure that covers the fixed true parameter in 95\% of repeated samples; a 95\% \omterm{def:qm:bayesian-methods:predictive}{credible
interval} contains the parameter with posterior probability 0.95 given this sample and the prior. They coincide numerically in simple
models with flat priors, and differ in meaning always.
\iqlookfor{Frequency over samples versus probability given data, and the role of the prior.}
\end{solution}

\begin{solution}{iq:qm:bayesian-methods:2}
Treat the fifty as draws from a population: estimate its mean and dispersion from the records (subtracting the noise variance), and shrink
each record toward the mean by $\sigma^2/(\sigma^2 + \tau^2)$. With three years and typical dispersion, two thirds of the best record's
excess is noise; expect about half its \omterm{def:qm:estimation:sharpe}{Sharpe ratio}.
\iqlookfor{Regression to the mean quantified, the noise variance $1/T$, and \omterm{def:qm:bayesian-methods:hier}{empirical Bayes}.}
\end{solution}

\begin{solution}{iq:qm:bayesian-methods:3}
Because total squared error over three or more means can be reduced by trading a little bias for a large cut in variance; pulling every
estimate toward a common point does that whatever the true means are. The catch: the gain is in the total, not for each unit, and an
unusual unit can be hurt.
\iqlookfor{Bias-variance trade in dimension three or more, and that dominance is for the sum.}
\end{solution}

\begin{solution}{iq:qm:bayesian-methods:4}
One never knows; one checks. Run several chains from dispersed starts and compare them ($\hat R$ near 1), look at trace plots, discard a
\omterm{def:qm:bayesian-methods:ess}{burn-in}, compute \omterm{def:qm:bayesian-methods:ess}{effective sample sizes} for every quantity reported, and compare with a different sampler or an analytic special case.
\iqlookfor{Multiple chains, $\hat R$, \omterm{def:qm:bayesian-methods:ess}{effective sample size}, and humility about convergence.}
\end{solution}

\begin{solution}{iq:qm:bayesian-methods:5}
Ridge regression's estimate is the posterior mode (and mean) of a linear model with a normal prior $\beta \sim \mathcal N(0, (\sigma^2/\lambda)I)$ on
the coefficients: the penalty $\lambda\lvert\beta\rvert^2$ is minus the log prior. The lasso is the mode under a Laplace prior.
\iqlookfor{Penalty as log prior, and the scale of $\lambda$ as a prior variance.}
\end{solution}

\begin{solution}{iq:qm:bayesian-methods:6}
Its acceptance rule makes the flow of probability from $x$ to $y$ equal to the flow from $y$ to $x$ under $\pi$ (\omterm{def:qm:markov-chains-and-queues:balance}{detailed balance}), so $\pi$ is
stationary, and an ergodic chain converges to it. A Gibbs step proposes from the exact conditional, which makes the Metropolis--Hastings
ratio exactly one.
\iqlookfor{\omterm{def:qm:markov-chains-and-queues:balance}{Detailed balance} written out, and the cancellation for Gibbs.}
\end{solution}
